\subsection{Other countries}\label{sec:countries}

This subsection shows that the two-floor structure is not a Danish peculiarity: it appears where the conditions it requires are met, and there too removing the housing allowance removes it.

\paragraph{Conditions.} The Danish and Norwegian results point to three conditions. (C1) A large share of the group lives on a statutory floor, so that its disposable income is fixed by statute. (C2) A second transfer is paid to a substantial part of that group but not to all of it. (C3) The second transfer separates the two parts: the distance between their medians is large relative to the spread of income within each. (C2) and (C3) decide whether the income of those on the floor is bimodal; (C1) decides whether that bimodality is visible in the group as a whole, as the Danish subsampling exercise of the previous subsection shows. In Denmark 71 to 76 per cent of the refined transfer-reliant group receive the housing benefit, and the distance between the medians of recipients and non-recipients is 2.8 to 3.5 times the interquartile range of non-recipients and 1.7 to 2.0 times that of recipients. Norway fails (C1), and (C3) as well.

(C2) can be quantified with the Danish data. Keep all non-recipients in the refined group of 2019 and a random subsample of recipients, so that recipients are 22 per cent of the group instead of 71. The test then rejects at the 5 per cent level in all ten subsamples, of about 7{,}300 people each, but in only two of ten when the subsample is cut to 1{,}000 people. With recipients at 10 per cent it rejects in none of ten subsamples at either size. A second transfer received by a fifth of the group can split it, but only in a large sample; one received by a tenth does not.

\paragraph{Data and design.} Most other LIS files are survey samples with a few hundred to a few thousand older people living alone in each wave, too few for the test (Appendix~\ref{app:sefi}). I therefore pool all waves from 2010 to 2021 for eight countries that record a housing allowance in LIS (item hi46) and have several thousand older people living alone across the waves: the Netherlands, the United Kingdom, Australia, Ireland, the United States, France, Sweden and Germany. Other LIS countries with the item were not examined. Pooling needs a common scale, because benefit rates and incomes change from wave to wave. In each wave income is expressed as a percentage of the median disposable income of transfer-reliant non-recipients of the allowance, the analogue of the Danish lower floor, so that this median is 100 in every wave. The group is the transfer-reliant group as defined above. The refinement that excludes public contributory pensions cannot be used outside Denmark, because in the United Kingdom, Ireland, the United States and Germany the main public old-age pension is itself recorded as contributory. Denmark is run through the same procedure, with its 2016 and 2019 waves, as a control. Table~\ref{tab:pool} reports the results. Because each wave is placed so that its non-recipients sit at 100, pooling cannot move a mode that no single wave has; it could, however, blur a second mode whose position varies from wave to wave, or produce one from waves whose single modes lie at different points. Two checks address this. Where the waves are large enough, the test is also run wave by wave; and for Denmark, the Netherlands, the United Kingdom, Australia and Ireland the pooled test is repeated with incomes divided by the median of all transfer-reliant people in the wave instead of non-recipients.

\begin{table}[t]
\centering\footnotesize
\caption{Transfer-reliant people aged 65 and over living alone (public pensions and benefits at least 90 per cent of income), all LIS waves 2010--2021 pooled. In each wave income is divided by the median disposable income of transfer-reliant non-recipients of a housing allowance and multiplied by 100, so that this median is 100 in every wave; household weights are normalised to average one within each wave. $n$: pooled number of transfer-reliant; share: of all aged 65 and over living alone; HB: share of the transfer-reliant receiving a housing allowance (LIS item hi46); benefit and gap: range over waves of the median allowance of recipients and of the difference between the median incomes of recipients and non-recipients, both on the same scale. $p$: bootstrap $p$-values of the test of unimodality, lowest 5 per cent trimmed, 199 draws, with cells of 2 and 4 points of the scale; less HB: income less the allowance, cells of 2; HB, no HB: recipients and non-recipients separately, cells of 2, tested when the pooled test rejects at the 10 per cent level. Denmark: waves 2016 and 2019 only, as a control; its no-HB and less-HB rejections come from former civil servants, and in the refined group of Table~\ref{tab:hb}, pooled in the same way, the $p$-values for less HB, HB and no HB are 1.000, 0.650 and 0.325.}\label{tab:pool}
\setlength{\tabcolsep}{3pt}
\begin{tabular}{lrrrrrrrrrrr}
\toprule
 & & & & & & & \multicolumn{5}{c}{$p$}\\
\cmidrule(lr){8-12}
Country & Waves & $n$ & Share & HB & Benefit & Gap & 2 & 4 & less HB & HB & no HB\\
\midrule
Denmark & 2 & 52{,}336 & 0.59 & 0.61 & 27--28 & 25 & 0.005 & 0.005 & 0.005 & 0.975 & 0.005\\
Netherlands & 12 & 3{,}315 & 0.16 & 0.65 & 19--23 & 18--25 & 0.010 & 0.005 & 0.195 & 1.000 & 0.005\\
United Kingdom & 12 & 13{,}185 & 0.38 & 0.66 & 41--47 & 47--55 & 0.005 & 0.005 & 0.315 & 0.005 & 0.065\\
Australia & 5 & 5{,}332 & 0.56 & 0.17 & 13--15 & 13--23 & 0.005 & 0.250 & 0.325 & 0.665 & 0.050\\
Ireland & 12 & 7{,}432 & 0.82 & 0.88 & 6--11 & $-40$ to 9 & 0.005 & 0.005 & 0.005 & 0.005 & 0.995\\
United States & 12 & 38{,}972 & 0.51 & 0.16 & 16--21 & $-7$ to 2 & 0.055 & 0.005 & 0.165 & 1.000 & 0.070\\
France & 12 & 8{,}964 & 0.10 & 0.87 & 10--22 & 3--41 & 0.190 & 0.450 & 0.150 & & \\
France 2010--14 & 5 & 3{,}939 & 0.10 & 0.89 & 18--22 & 38--41 & 0.005 & 0.005 & 0.425 & & \\
Sweden & 12 & 2{,}640 & 0.24 & 0.63 & 22--30 & 0--17 & 1.000 & 0.855 & 0.125 & & \\
Germany & 12 & 10{,}923 & 0.66 & 0.04 & 4--8 & $-45$ to $-33$ & 0.685 & 0.645 & 0.925 & & \\
\bottomrule
\end{tabular}
\end{table}


\paragraph{Replications.} The procedure reproduces the Danish result: the pooled Danish transfer-reliant group is bimodal ($p=0.005$) with modes at 96 and 126 on the scale, as is each of the two waves and the pool on the alternative scale, and in the refined group recipients, non-recipients and income less the benefit are each unimodal.

The Netherlands has the Danish structure on a smaller scale. The flat-rate state pension (AOW) is the floor for 16 per cent of older people living alone; 65 per cent of them receive the housing allowance (huurtoeslag), whose median is 19 to 23 per cent of the floor. The pooled test rejects ($p=0.010$ with cells of 2 points, $0.005$ with cells of 4), with modes at 100 and 124: the floor, and the floor plus the allowance. Recipients alone are unimodal ($p=1$), and subtracting the allowance from income removes the valley ($p=0.195$). On the alternative scale the test also rejects ($p=0.005$). The Dutch waves have 120 to 480 transfer-reliant people each, too few to test one at a time, but in every wave the cells of the pooled valley hold less mass than the cells of equal width around the upper mode. Non-recipients show a further mode at 60 per cent of the floor. This is consistent with the rule that the AOW accrues at 2 per cent for each year of residence, so that people who spent part of their working lives abroad receive a partial pension; I do not pursue it.

In the United Kingdom the transfer-reliant are 38 per cent of older people living alone, and two thirds of them receive support with housing costs, which LIS records in hi46 and which comprises Housing Benefit and council-tax support. The pooled distribution is bimodal ($p=0.005$ at both cell widths), with modes at 102 and 148, and it is unimodal once the housing support is subtracted ($p=0.315$). On the alternative scale the test rejects as well ($p=0.005$). Wave by wave, with 1{,}000 to 1{,}400 people each (524 in 2020), it rejects at the 5 per cent level in 2011, 2016 and 2021, and in these waves the modes lie at 100 to 104 and 136 to 152, where the pooled modes are. Unlike the Danish recipients, the British recipients are themselves bimodal. Housing Benefit pays a share of the actual rent, not a sum bounded by a low maximum, and council-tax support on its own is small, so recipients do not share one amount. When recipients are split by the size of the payment, measured here against the median income of all transfer-reliant in the wave, recipients of more than 15 per cent of that median are unimodal ($p=0.99$), and without the fifth of the group who receive smaller amounts the distribution is only weakly bimodal ($p=0.07$). The British structure is thus the floor, on which non-recipients and recipients of council-tax support alone both stand, and the floor plus Housing Benefit.

France shows the structure appearing and disappearing with (C3). Ten per cent of older people living alone are transfer-reliant, and nine in ten of them receive a housing allowance. From 2010 to 2014 the median income of recipients was 38 to 41 per cent above that of non-recipients, and the pooled distribution of these five waves is bimodal ($p=0.005$ at both cell widths) with modes at 100 and 132; subtracting the allowance removes the valley ($p=0.425$). The gap then shrank, to 3 per cent in 2021, and with all twelve waves pooled the second mode is smeared out ($p=0.19$).

Australia is the marginal case that (C2) predicts, and the evidence there is marginal too. Fifty-six per cent of older people living alone live on the Age Pension, but only 17 per cent of them receive Rent Assistance, which amounts to 13 to 15 per cent of the floor. With cells of 2 points the pooled test rejects ($p=0.005$), with modes at 100 and 114, and subtracting Rent Assistance removes the valley ($p=0.325$). But with cells of 4 points, which make the gap between the modes three cells wide, it does not reject ($p=0.25$); on the alternative scale it gives $p=0.065$; and no single wave rejects at the 5 per cent level. Non-recipients alone give $p=0.05$, with a shallow dip just below the floor (modes at 90 and 100); whatever produces it, it is not Rent Assistance, and it lies below the valley at 109 to 113 that the pooled test finds. The Danish experiment says that a recipient share of this size is detectable only in large samples, so the Australian result is what the conditions predict, but it is not a replication.

\paragraph{Countries without the structure.} In Germany only 4 per cent of the transfer-reliant receive a housing allowance, and recipients are poorer, not richer, than non-recipients. (C2) fails, and the test does not reject. In Sweden recipients of the housing supplement are barely richer than non-recipients: the gap between the medians is 0 to 17 points of the scale, while the interquartile range among non-recipients is 26 to 36 points. (C3) fails, and the test does not reject.  In the United States housing assistance reaches 16 per cent of the transfer-reliant and does not separate them from the rest (the gap is $-7$ to 2 points), and non-recipients are widely spread (interquartile range 57 to 72 points), so there is no floor for the allowance to lift anyone above: (C3) fails. With cells of 4 points the test finds a shallow second mode 8 points above the first ($p=0.005$, excess mass 0.004 against 0.026 in Denmark). It is unrelated to housing assistance, since recipients alone give $p=1$; the two cell widths locate the dip in different places, and at this sample size the test is slightly liberal (Section~\ref{sec:data}). I read it as noise. In Ireland 88 per cent of the transfer-reliant receive a housing-related payment, and it is small (6 to 11 points), so (C3) fails. The Irish pooled test rejects, but the rejection survives the subtraction of the payment ($p=0.005$), and its location depends on the scale: when incomes are divided by the median of all transfer-reliant rather than of non-recipients, the valley is a single cell 2 points wide, with the two modes in the adjacent cells at 96 and 100. Whatever produces it, it is not the housing payment.

The conditions thus sort the countries. Where a housing allowance is paid to a substantial part of a group living on a flat floor and lifts them clearly above it (Denmark, the Netherlands, the United Kingdom, and France until 2014), the income of the group has two modes, and subtracting the allowance removes the second. Where the recipient share is small (Australia) the evidence is weak, as the Danish experiment says it must be. Where a condition fails (Norway, Sweden, Germany, the United States, Ireland, and France after the gap closed), no second mode is due to the allowance.